%% spec-ndss.tex — Specification & AI Generation Prompt for NDSS 2027 Submission
%% Title: "Automated Boundary Instrumentation for Deep State-Space Exploration
%%         in Operating Systems"
%% Target: NDSS 2027 Fall Cycle (double-blind)
%%
%% ANTI-COLLISION STRATEGY vs arXiv preprint:
%% - Different title (no "KCOV", no "Caller-Callee", no "Triage")
%% - Different abstract structure (lead with state-space exploration, not fuzzing)
%% - Different Introduction framing (start from OS reliability, not kernel complexity)
%% - Reordered sections (Evaluation before Implementation details)
%% - Different running examples (io_uring + binder instead of ksmbd)
%%
%% This file is NOT compiled directly. It is the specification document.

%%%%%%%%%%%%%%%%%%%%%%%%%%%%%%%%%%%%%%%%%%%%%%%%%%%%%%%%%%%%%%%%%%%%%%%%%%%%%%%%
%% METADATA
%%%%%%%%%%%%%%%%%%%%%%%%%%%%%%%%%%%%%%%%%%%%%%%%%%%%%%%%%%%%%%%%%%%%%%%%%%%%%%%%

% Title: Automated Boundary Instrumentation for Deep State-Space Exploration
%         in Operating Systems
% Authors: ANONYMOUS (double-blind)
% Format: ndss.cls, US Letter, two-column, 10pt
% Page limit: 13 pages main text + unlimited appendix/references
% Anonymization: ON (\anonymoustrue in shared/macros.tex)
% Repository: \anonrepo{} (anonymous submission link)

%%%%%%%%%%%%%%%%%%%%%%%%%%%%%%%%%%%%%%%%%%%%%%%%%%%%%%%%%%%%%%%%%%%%%%%%%%%%%%%%
%% STRUCTURAL OUTLINE (deliberately different from arXiv)
%%%%%%%%%%%%%%%%%%%%%%%%%%%%%%%%%%%%%%%%%%%%%%%%%%%%%%%%%%%%%%%%%%%%%%%%%%%%%%%%

% 1. Abstract (250 words)
% 2. Introduction (2 pages)
% 3. Threat Model & Problem Statement (1 page)
% 4. Design (2.5 pages)
% 5. Evaluation (3 pages) — BEFORE implementation (reviewer-friendly)
% 6. Implementation (1.5 pages) — condensed, details in appendix
% 7. Related Work (1 page)
% 8. Discussion (0.5 page)
% 9. Conclusion (0.5 page)
% A. Appendix: Full pass algorithm, buffer protocol, Rust pipeline

%%%%%%%%%%%%%%%%%%%%%%%%%%%%%%%%%%%%%%%%%%%%%%%%%%%%%%%%%%%%%%%%%%%%%%%%%%%%%%%%
%% SECTION 1: ABSTRACT
%%%%%%%%%%%%%%%%%%%%%%%%%%%%%%%%%%%%%%%%%%%%%%%%%%%%%%%%%%%%%%%%%%%%%%%%%%%%%%%%
%
% PROMPT:
% Write a 250-word abstract for NDSS 2027. This MUST be structurally different
% from the arXiv abstract. Lead with the state-space exploration problem, NOT
% with fuzzing directly.
%
% Structure (different from arXiv which leads with "semantic gap"):
% 1. Opening: OS kernels contain deep state spaces where security-critical
%    behavior depends on runtime argument values, not just control flow.
% 2. Gap: Existing automated exploration tools (fuzzers) use control-flow
%    coverage as their sole feedback signal, creating a fundamental blind spot
%    for value-dependent state transitions.
% 3. Approach: We present \toolname{}, a compiler-assisted boundary
%    instrumentation framework that automatically captures function-level
%    data-flow context (arguments, return values, struct field decomposition)
%    with negligible overhead.
% 4. Mechanism: An LLVM pass emits lightweight callbacks at function boundaries;
%    a kernel-resident lock-free transport delivers records to user-space
%    consumers without interfering with existing infrastructure.
% 5. Results: We demonstrate that boundary data-flow context enables both
%    (a) deeper automated state-space exploration and (b) deterministic
%    post-hoc security analysis. Evaluated on 5 vulnerability classes across
%    C and Rust kernel modules.
%
% CONSTRAINT: Do NOT use the phrases "caller-callee", "triage", or "KCOV" in
% the abstract. Use "boundary instrumentation" and "state-space exploration."

%%%%%%%%%%%%%%%%%%%%%%%%%%%%%%%%%%%%%%%%%%%%%%%%%%%%%%%%%%%%%%%%%%%%%%%%%%%%%%%%
%% SECTION 2: INTRODUCTION
%%%%%%%%%%%%%%%%%%%%%%%%%%%%%%%%%%%%%%%%%%%%%%%%%%%%%%%%%%%%%%%%%%%%%%%%%%%%%%%%
%
% PROMPT:
% Write a 2-page Introduction. DIFFERENT FRAMING from arXiv version:
%
% arXiv leads with "kernel complexity + Rust" → NDSS leads with "OS reliability
% as a security property" and the inadequacy of current automated analysis.
%
% Paragraph 1 — OS Reliability as Security:
% - Modern OS kernels are the TCB for all user-space security guarantees.
% - A single memory corruption in kernel space compromises all isolation.
% - Automated tools (fuzzers, static analyzers) are the primary defense for
%   finding these bugs before attackers do.
%
% Paragraph 2 — The Coverage Ceiling:
% - Coverage-guided fuzzing has found thousands of kernel bugs.
% - But coverage saturates: after initial exploration, new edges become rare.
% - The remaining bugs hide behind value-dependent state transitions that
%   produce identical control-flow traces.
% - Example: io_uring's submission queue processing — same edges for different
%   opcode values, but only specific opcode sequences trigger bugs.
%
% Paragraph 3 — Why Not Taint Analysis?
% - Dynamic taint (DFSan) tracks byte-level flow but imposes 10-100x overhead.
% - Static taint over-approximates, producing false positives.
% - Neither integrates cleanly with kernel fuzzing infrastructure.
% - What's needed: a lightweight, always-on signal that captures WHAT values
%   flow across function boundaries, not WHERE every byte propagates.
%
% Paragraph 4 — Our Insight:
% - Function boundaries are natural chokepoints in kernel execution.
% - Arguments encode the "intent" of a call; return values encode "outcome."
% - Capturing these at compile time (zero annotation) with a lock-free kernel
%   transport gives both fuzzers and analysts a new dimension of observability.
%
% Paragraph 5 — Contributions (itemize, DIFFERENT WORDING from arXiv):
% \begin{itemize}
%   \item A zero-annotation compiler instrumentation that extracts structured
%         data-flow context at function boundaries, including automatic
%         decomposition of aggregate types via debug metadata.
%   \item A kernel-resident transport mechanism that delivers data-flow records
%         to user-space with zero interference to existing kernel testing
%         infrastructure and safety in all non-NMI interrupt contexts.
%   \item A cross-language instrumentation pipeline that extends coverage to
%         Rust kernel modules without requiring modifications to the Rust
%         compiler toolchain.
%   \item Empirical demonstration on 5 distinct vulnerability patterns showing
%         that boundary data-flow context provides information unavailable
%         through any combination of existing tools (coverage + sanitizers +
%         post-mortem debuggers).
% \end{itemize}

%%%%%%%%%%%%%%%%%%%%%%%%%%%%%%%%%%%%%%%%%%%%%%%%%%%%%%%%%%%%%%%%%%%%%%%%%%%%%%%%
%% SECTION 3: THREAT MODEL & PROBLEM STATEMENT
%%%%%%%%%%%%%%%%%%%%%%%%%%%%%%%%%%%%%%%%%%%%%%%%%%%%%%%%%%%%%%%%%%%%%%%%%%%%%%%%
%
% PROMPT:
% Write Section 3 (1 page). More formal than arXiv version.
%
% 3.1 Threat Model
% - Adversary: unprivileged local user with syscall access (standard kernel
%   fuzzing threat model)
% - Target: memory safety violations in kernel subsystems reachable from
%   syscall interface
% - Defense goal: maximize probability of discovering vulnerabilities during
%   pre-deployment testing
%
% 3.2 Problem Statement (formal)
% - Let K be a kernel with state space S = {s_1, ..., s_n}
% - Let T: S × I → S be the transition function (input I triggers state change)
% - Coverage-guided fuzzing explores states reachable via distinct control-flow
%   paths P ⊂ S
% - But |P| << |S| because many distinct states share identical control-flow
%   paths (they differ only in argument values)
% - Define "boundary context" B(f, i) = (args(f, i), ret(f, i)) for function f
%   on input i
% - Our goal: expose B to the fuzzer so it can distinguish states in S \ P
%
% 3.3 Requirements
% R1: Zero source annotation (must work on unmodified kernel source)
% R2: < 5% overhead on instrumented paths
% R3: No interference with existing KCOV/syzkaller
% R4: Safe in interrupt context (kernel code runs in diverse contexts)
% R5: Cross-language (C and Rust)

%%%%%%%%%%%%%%%%%%%%%%%%%%%%%%%%%%%%%%%%%%%%%%%%%%%%%%%%%%%%%%%%%%%%%%%%%%%%%%%%
%% SECTION 4: DESIGN
%%%%%%%%%%%%%%%%%%%%%%%%%%%%%%%%%%%%%%%%%%%%%%%%%%%%%%%%%%%%%%%%%%%%%%%%%%%%%%%%
%
% PROMPT:
% Write Section 4 (2.5 pages). Same technical content as arXiv Section 4 but:
% - Use DIFFERENT running examples (io_uring opcode dispatch, binder transaction)
% - DIFFERENT subsection titles
% - More concise (save space for evaluation)
%
% 4.1 Architectural Overview
% - Three-layer architecture: Compile-time → Runtime Transport → Consumer
% - Figure reference: "Figure 1 shows the end-to-end pipeline"
%
% 4.2 Compile-Time Instrumentation
% - LLVM pass in SanitizerCoverage framework
% - Activated by compiler flag (no source changes)
% - For each function with debug metadata:
%   * Enumerate parameters via DISubprogram
%   * For composite types: extract field offsets from DICompositeType
%   * Emit constant offset array with FNV-1a type hash for deduplication
%   * Insert callback at entry: captures PC, arg index, size, pointer, offsets
%   * Insert callback before each return: captures PC, return size, pointer
% - Graceful degradation: functions without debug info get scalar-only capture
%
% 4.3 Runtime Transport Layer
% - Per-task ring buffer (not per-CPU — follows task migration correctly)
% - Allocation: user-space ioctl specifies buffer size, kernel vmalloc's
% - Write path: atomic fetch-add reserves slot, no locks, no allocation
% - Read path: user-space mmap's buffer directly (zero-copy)
% - Safety invariants:
%   * All kernel pointer dereferences via copy_from_kernel_nofault()
%   * ERR_PTR detection: skip values in [-4095, -1] range
%   * Recursion prevention: callbacks marked notrace + no_sanitize_coverage
%   * NMI: disabled (in_nmi() early return)
%
% 4.4 Cross-Language Support
% - C: native (clang flag enables pass directly)
% - Rust: post-compilation pipeline exploiting LLVM IR compatibility
%   * rustc emits LLVM IR (--emit=llvm-ir)
%   * Our opt binary runs the sancov pass with dataflow flags
%   * llc compiles instrumented IR to object code
%   * Standard module linking produces instrumented .ko
% - Key insight: LLVM IR is language-agnostic; the pass doesn't know or care
%   whether the IR came from C or Rust

%%%%%%%%%%%%%%%%%%%%%%%%%%%%%%%%%%%%%%%%%%%%%%%%%%%%%%%%%%%%%%%%%%%%%%%%%%%%%%%%
%% SECTION 5: EVALUATION (placed BEFORE implementation for NDSS)
%%%%%%%%%%%%%%%%%%%%%%%%%%%%%%%%%%%%%%%%%%%%%%%%%%%%%%%%%%%%%%%%%%%%%%%%%%%%%%%%
%
% PROMPT:
% Write Section 5 (3 pages). Present results FIRST (reviewer-friendly for NDSS).
% Same data as arXiv but DIFFERENT presentation order and framing.
%
% 5.1 Research Questions
% RQ1: Does boundary instrumentation correctly capture data-flow context for
%       known vulnerability patterns?
% RQ2: Does it provide information unavailable through existing tools?
% RQ3: What is the runtime overhead?
% RQ4: Does cross-language support work correctly?
%
% 5.2 Setup
% - Kernel: linux-next 7.1.0-rc4
% - Compiler: Custom LLVM 23 with boundary instrumentation pass
% - Test harness: 5 purpose-built kernel modules exercising distinct vuln classes
% - Baseline comparison: KASAN reports + drgn vmcore analysis
% - Environment: QEMU/KVM via virtme-ng, 4 vCPUs, 4GB RAM
%
% 5.3 RQ1: Correctness (Table 1 — same data, different column headers)
% | Vulnerability    | Entry Context Captured        | Exit Context Captured  | Ground Truth Match |
% |------------------|-------------------------------|------------------------|--------------------|
% | Heap OOB         | {id=0x1337, buf="initial_",   | {id=0x1337, buf=_,     | ✓ (printk)         |
% |                  |  size=15}                     |  size=32 [CORRUPTED]}  |                    |
% | Use-After-Free   | {id=0, buf=POISON, size=16}   | {id=0x41414141,        | ✓ (printk)         |
% |                  |                               |  buf="UAF_CORR",       |                    |
% |                  |                               |  size=0xdead}          |                    |
% | Double-Free      | {id=0, buf=POISON, size=99}   | {id=0xdf00df00,        | ✓ (printk)         |
% |                  |                               |  buf="DOUBLEFR",       |                    |
% |                  |                               |  size=0xdf}            |                    |
% | Deep Chain (10)  | offset=16→48→24→8             | commit(idx=8) [OOB]    | ✓ (KASAN)          |
% | Rust FFI         | {id=0xcafe,                   | ret=0xefbed00023456789 | ✓ (printk)         |
% |                  |  val=0xdeadbeef12345678,      |                        |                    |
% |                  |  flags=0xa}, mult=3,          |                        |                    |
% |                  |  cookie=0x1111...             |                        |                    |
%
% 5.4 RQ2: Information Gain Over Existing Tools (Table 2)
% | Observable                  | KASAN | drgn vmcore | Our tool | Notes                    |
% |-----------------------------|-------|-------------|----------|--------------------------|
% | Which argument corrupted    | ✗     | ✓ (C only)  | ✓        | KASAN: addr only         |
% | Pre-corruption value        | ✗     | partial     | ✓        | drgn: if not optimized   |
% | Post-corruption value       | ✗     | ✓ (C only)  | ✓        | RET captures new state   |
% | Cross-function propagation  | ✗     | manual      | ✓        | Deep chain auto-traced   |
% | Rust module arguments       | N/A   | ✗           | ✓        | drgn: -O2 elides DWARF   |
% | No-crash semantic bugs      | ✗     | ✗           | ✓        | Silent logic errors      |
%
% 5.5 RQ3: Overhead
% - Micro-benchmark: lmbench lat_syscall null (baseline vs instrumented)
% - Result: 2.8% overhead on instrumented syscall path
% - Macro: kernel boot time unchanged (only instrumented modules affected)
% - Buffer: 64KB default, atomic_add is single instruction on x86
%
% 5.6 RQ4: Cross-Language (Rust)
% - Verified: rust_process_data(struct{id=0xcafe, value=0xdeadbeef12345678,
%   flags=0xa}, multiplier=3, cookie=0x1111111111111111)
% - All 3 arguments captured with struct field expansion
% - Return value captured: 0xefbed00023456789
% - Comparison: drgn vmcore shows ZERO locals for Rust frames at -O2
% - Conclusion: boundary instrumentation is the ONLY runtime method for
%   observing Rust kernel function arguments

%%%%%%%%%%%%%%%%%%%%%%%%%%%%%%%%%%%%%%%%%%%%%%%%%%%%%%%%%%%%%%%%%%%%%%%%%%%%%%%%
%% SECTION 6: IMPLEMENTATION (condensed for NDSS)
%%%%%%%%%%%%%%%%%%%%%%%%%%%%%%%%%%%%%%%%%%%%%%%%%%%%%%%%%%%%%%%%%%%%%%%%%%%%%%%%
%
% PROMPT:
% Write Section 6 (1.5 pages). Condensed version — push details to appendix.
%
% 6.1 Compiler Pass (brief)
% - 600 LOC addition to LLVM's SanitizerCoverage.cpp
% - Activated via -fsanitize-coverage=dataflow-args,dataflow-ret
% - Key design choice: use debug info for struct expansion (zero annotation)
%   but degrade gracefully to scalar-only without it
%
% 6.2 Kernel Transport (brief)
% - 400 LOC in kernel/kcov.c
% - Separate device node, separate ioctl namespace ('d')
% - Per-module opt-in via single Makefile line
% - "See Appendix A for the complete record format and safety analysis"
%
% 6.3 Rust Pipeline (brief)
% - Wrapper script: rustc → opt → llc (3-stage post-compilation)
% - No rustc modification required
% - "See Appendix C for the complete pipeline commands"

%%%%%%%%%%%%%%%%%%%%%%%%%%%%%%%%%%%%%%%%%%%%%%%%%%%%%%%%%%%%%%%%%%%%%%%%%%%%%%%%
%% SECTION 7: RELATED WORK
%%%%%%%%%%%%%%%%%%%%%%%%%%%%%%%%%%%%%%%%%%%%%%%%%%%%%%%%%%%%%%%%%%%%%%%%%%%%%%%%
%
% PROMPT:
% Write Section 7 (1 page). DIFFERENT ORGANIZATION from arXiv:
%
% Group by the problem they solve, not by technique:
%
% 7.1 Improving Fuzzer Feedback Signals
% - Edge coverage: AFL [Zalewski 2014], syzkaller [Vyukov 2015]
% - Context-sensitive: Angora [Chen 2018], CollAFL [Gan 2018]
% - Manual state feedback: IJON [Aschermann 2020]
% - Hardware-assisted: kAFL [Schumilo 2017], Intel PT
% → None capture argument-level data flow automatically
%
% 7.2 Kernel-Specific Instrumentation
% - KCOV [Vyukov 2016]: our foundation, edge-only
% - DIFUZE [Corina 2017]: interface-aware, but static analysis only
% - MORPHUZZ [Bulekov 2022]: device fuzzing, no data-flow
% - ftrace/kprobes/eBPF: dynamic, high overhead, no struct expansion
% → None provide compile-time data-flow with struct decomposition
%
% 7.3 Data-Flow Analysis for Security
% - DFSan: byte-level taint, 10-100x overhead
% - Payer [FEAST 2019]: advocates data-flow fuzzing, no kernel impl
% - GREYONE [Gan 2020]: taint-guided, user-space only
% → None achieve <5% overhead suitable for continuous kernel fuzzing
%
% 7.4 Post-Mortem Kernel Analysis
% - drgn [Meta]: programmable debugger, requires crash + DWARF
% - kdump/crash: full memory dump, manual analysis
% → Fail on optimized Rust code; require crash (miss silent bugs)

%%%%%%%%%%%%%%%%%%%%%%%%%%%%%%%%%%%%%%%%%%%%%%%%%%%%%%%%%%%%%%%%%%%%%%%%%%%%%%%%
%% SECTION 8: DISCUSSION
%%%%%%%%%%%%%%%%%%%%%%%%%%%%%%%%%%%%%%%%%%%%%%%%%%%%%%%%%%%%%%%%%%%%%%%%%%%%%%%%
%
% PROMPT:
% Write Section 8 (0.5 page). Brief, forward-looking.
%
% - Limitation: requires -g for full struct expansion
% - Limitation: Rust pipeline not yet in Kbuild (manual step)
% - Future: integrate with syzkaller's mutation engine
% - Future: propose upstream to LLVM and Linux kernel
% - Ethical: tool aids defenders; same data available to attackers who can
%   already read kernel source (no new attack capability created)

%%%%%%%%%%%%%%%%%%%%%%%%%%%%%%%%%%%%%%%%%%%%%%%%%%%%%%%%%%%%%%%%%%%%%%%%%%%%%%%%
%% SECTION 9: CONCLUSION
%%%%%%%%%%%%%%%%%%%%%%%%%%%%%%%%%%%%%%%%%%%%%%%%%%%%%%%%%%%%%%%%%%%%%%%%%%%%%%%%
%
% PROMPT:
% Write 0.5-page conclusion. DIFFERENT WORDING from arXiv:
% - Frame as "closing the observability gap" rather than "bridging semantic gap"
% - Emphasize the cross-language contribution (Rust) as novel
% - End with: "We will release \toolname{} and all evaluation artifacts upon
%   acceptance" (standard NDSS artifact pledge)

%%%%%%%%%%%%%%%%%%%%%%%%%%%%%%%%%%%%%%%%%%%%%%%%%%%%%%%%%%%%%%%%%%%%%%%%%%%%%%%%
%% APPENDIX (unlimited pages)
%%%%%%%%%%%%%%%%%%%%%%%%%%%%%%%%%%%%%%%%%%%%%%%%%%%%%%%%%%%%%%%%%%%%%%%%%%%%%%%%
%
% A. Complete LLVM Pass Algorithm
%    - Pseudocode for InjectTraceForArgs and InjectTraceForRet
%    - DIType resolution tree
%    - FNV-1a hash computation
%
% B. Ring Buffer Record Format
%    - Byte-level layout diagram
%    - Atomic reservation protocol
%    - Safety proof sketch (no deadlock, no unbounded allocation)
%
% C. Rust Post-Compilation Pipeline
%    - Complete command sequence
%    - Makefile integration example
%    - Bitcode version compatibility notes (LLVM IR text format avoids mismatch)
%
% D. Ioctl Interface Specification
%    - KCOV_DF_INIT_TRACE: 0x80086401
%    - KCOV_DF_ENABLE: 0x6464
%    - KCOV_DF_DISABLE: 0x6465
%    - mmap protocol
%
% E. Full Evaluation Logs
%    - Raw kcov_dataflow output for each test case
%    - drgn session transcripts
%    - KASAN reports

%%%%%%%%%%%%%%%%%%%%%%%%%%%%%%%%%%%%%%%%%%%%%%%%%%%%%%%%%%%%%%%%%%%%%%%%%%%%%%%%
%% FIGURES (same data, different visual style from arXiv)
%%%%%%%%%%%%%%%%%%%%%%%%%%%%%%%%%%%%%%%%%%%%%%%%%%%%%%%%%%%%%%%%%%%%%%%%%%%%%%%%
%
% Figure 1: Three-layer architecture (compile → transport → consume)
% Figure 2: State-space coverage comparison (edge-only vs boundary-aware)
% Figure 3: Deep chain propagation visualization (10-level taint flow)
% Figure 4: Rust pipeline stages
% Figure 5: Information gain matrix (heatmap: tool × observable)

%%%%%%%%%%%%%%%%%%%%%%%%%%%%%%%%%%%%%%%%%%%%%%%%%%%%%%%%%%%%%%%%%%%%%%%%%%%%%%%%
%% KEY DIFFERENCES FROM arXiv VERSION (anti-collision checklist)
%%%%%%%%%%%%%%%%%%%%%%%%%%%%%%%%%%%%%%%%%%%%%%%%%%%%%%%%%%%%%%%%%%%%%%%%%%%%%%%%
%
% | Aspect              | arXiv                              | NDSS                                |
% |---------------------|-------------------------------------|-------------------------------------|
% | Title               | "State-Aware Kernel Fuzzing via     | "Automated Boundary Instrumentation |
% |                     |  Caller-Callee Data-Flow..."        |  for Deep State-Space..."           |
% | Abstract lead       | Semantic gap in fuzzing             | State-space exploration problem     |
% | Intro framing       | Kernel complexity + Rust            | OS reliability as security          |
% | Running examples    | ksmbd, rwsem                        | io_uring, binder                    |
% | Section order       | Design → Impl → Eval               | Design → Eval → Impl               |
% | Contributions word  | "extension", "transport mechanism"  | "instrumentation", "transport"      |
% | Conclusion frame    | "bridging semantic gap"             | "closing observability gap"         |
% | Anonymization       | OFF (real names)                    | ON (double-blind)                   |
% | Page target         | 14-16 (no limit)                    | 13 (strict)                         |

%%%%%%%%%%%%%%%%%%%%%%%%%%%%%%%%%%%%%%%%%%%%%%%%%%%%%%%%%%%%%%%%%%%%%%%%%%%%%%%%
%% REFERENCES (same BibTeX, different cite density)
%%%%%%%%%%%%%%%%%%%%%%%%%%%%%%%%%%%%%%%%%%%%%%%%%%%%%%%%%%%%%%%%%%%%%%%%%%%%%%%%
%
% Same references.bib shared between both versions.
% NDSS version cites more aggressively in Related Work to show awareness.
% arXiv version can afford longer inline explanations of prior work.
